% from ega1-4-fr.tex
\oldpage[I]{127}
un point $x\in X$ s'il existe un voisinage ouvert $U$ de $x$ dans $X$ tel
que la restriction de $f$ au préschéma induit $U$ soit une immersion ouverte
de $U$ dans $Y$. On dit que $f$ est un isomorphisme local si $f$ est un
isomorphisme local en tout point de $X$.
\end{definition}

\begin{env}[4.5.3]
\label{I.4.5.3-fr}
Une immersion (resp. une immersion fermée) $f:X\to Y$ peut donc se
caractériser comme une immersion locale telle que $f$ soit un homéomorphisme
de l'espace sous-jacent à $X$ sur une partie (resp. une partie fermée) de
$Y$. Une immersion ouverte $f$ peut se caractériser comme un isomorphisme
local \emph{injectif}.
\end{env}

\begin{proposition}[4.5.4]
\label{I.4.5.4-fr}
Soient $X$ un préschéma irréductible, $f:X\to Y$ un morphisme injectif
dominant. Si $f$ est une immersion locale, $f$ est une immersion et $f(X)$
est ouvert dans $Y$.
\end{proposition}

En effet, soit $x\in X$ et soient $U$ un voisinage ouvert de $x$, $V$ un
voisinage ouvert de $f(x)$ dans $Y$ tels que la restriction de $f$ à $U$
soit une immersion fermée dans $V$~; comme $U$ est dense dans $X$, $f(U)$
est dense dans $Y$ par hypothèse, donc $f(U)=V$ et $f$ est un homéomorphisme
de $U$ sur $V$~; l'hypothèse que $f$ est injectif entraîne que
$f^{-1}(V)=U$, d'où aussitôt la proposition.

\begin{proposition}[4.5.5]
\label{I.4.5.5-fr}
\begin{enumerate}
  \item[\rm(i)] Le composé de deux immersions locales (resp. de deux
    isomorphismes locaux) est une immersion locale (resp. un isomorphisme
    local).
  \item[\rm(ii)] Soient $f:X\to X'$, $g:Y\to Y'$ deux $S$-morphismes. Si
    $f$ et $g$ sont des immersions locales (resp. des isomorphismes locaux),
    il en est de même de $f\times_S g$.
  \item[\rm(iii)] Si un $S$-morphisme $f$ est une immersion locale (resp. un
    isomorphisme local), il en est de même de $f_{(S')}$ pour toute extension
    $S'\to S$ du préschéma de base.
\end{enumerate}
\end{proposition}

D'après (\hyperref[I.3.5.1-fr]{3.5.1}), il suffit de prouver (i) et (ii).

(i) résulte aussitôt de la transitivité des immersions fermées (resp.
ouvertes) (\hyperref[I.4.2.4-fr]{4.2.4}) et du fait que si $f$ est un
homéomorphisme de $X$ sur une partie fermée de $Y$, pour tout ouvert
$U\subset X$, $f(U)$ est ouvert dans $f(X)$, donc il existe une partie
ouverte $V$ de $Y$ telle que $f(U)=V\cap f(X)$, et $f(U)$ est par suite fermé
dans $V$.

Pour démontrer (ii), soient $p$, $q$ les projections de $X\times_S Y$, $p'$,
$q'$ celles de $X'\times_S Y'$. Il existe par hypothèse des voisinages
ouverts $U$, $U'$, $V$, $V'$ de $x=p(z)$, $x'=p'(z')$, $y=q(z)$,
$y'=q'(z')$ respectivement, tels que les restrictions de $f$ et $g$ à $U$
et $V$ respectivement soient des immersions fermées (resp. ouvertes) dans
$U'$ et $V'$ respectivement. Comme l'espace sous-jacent de $U\times_S V$ et
celui de $U'\times_S V'$ s'identifient aux voisinages ouverts
$p^{-1}(U)\cap q^{-1}(V)$ et ${p'}^{-1}(U')\cap{q'}^{-1}(V')$ de $z$ et
$z'$ respectivement (\hyperref[I.3.2.7-fr]{3.2.7}), la proposition résulte
de (\hyperref[I.4.3.1-fr]{4.3.1}).
